\section{The Target: \FCdot}\label{sec:target}

\begin{figure*}
\small
\[
\begin{array}{@{}lrcl@{}}
\text{Types} & S,T & ::= & \bot \mid \sel{x}{\ell} \mid \pit{S}{T} \mid \obj{\tel} \qquad\quad \top \mathrel{:=} \obj{\tnil} \\
\text{Propositions} & P & ::= & \ple{S}{T} \mid \peq{S}{T} \mid \phas{\ell} \\
\text{Telescopes} & \tel & ::= & \tnil \mid \tcons{\tel}{P} \\[2pt]
\text{Inclusion} & e & ::= & \ev{refl}\,T \mid \ev{trans}\,e\,e \mid \ev{top}\,T \mid \ev{bot}\,T \mid \ev{eqToLe}\,\varphi \mid \ev{pi}\,e\,e \\
& & \mid & \ev{obj}\,\tel\,m \mid \ev{pair}\,\tel\,\tel\,e\,e \mid \ev{member}\,a\,e\,i \\
\text{Equality} & \varphi & ::= & \ev{refl}\,T \mid \ev{symm}\,\varphi \mid \ev{trans}\,\varphi\,\varphi \mid \ev{def}\,x\,\ell \mid \ev{member}\,a\,e\,i \\
\text{Presence} & h & ::= & \ev{member}\,a\,e\,i \mid \ev{field}\,\ell \\
\text{Holes} & H & ::= & \ev{le}\,j \mid \ev{eq}\,j \mid \ev{eqSym}\,j \\
\text{Sides} & q & ::= & \ev{none} \mid \ev{some}\,e \\
\text{Morphisms} & m & ::= & \ev{nil} \mid \ev{le}\,m\,q\,H\,q \mid \ev{eq}\,m\,j\,b \mid \ev{has}\,m\,j \\[2pt]
\text{Atoms} & a,b & ::= & x \mid \ev{cast}\,a\,e \mid \ev{fold}\,\tel\,a \mid \ev{unfold}\,a \mid \ev{both}\,\tel\,\tel\,a\,a \\
\text{Terms} & t,u & ::= & a \mid v \mid a\,b \mid \ev{proj}\,a\,\ell\,h \mid \klet x = t \kin u \mid \ev{cast}\,t\,e \\
\text{Values} & v & ::= & \lambda(x : S).\,t \mid \lit{W}{F} \mid \ev{cast}\,v\,e \\
\text{Witnesses} & W & ::= & \tnil \mid W, \ell = T \qquad \text{Fields} \quad F ::= \tnil \mid F, \ell = t \\
\text{Contexts} & \Gamma & ::= & \cdot \mid \Gamma, x : T \mid \Gamma, x : T \mathrel{:=} W ; \vec{\ell}
\end{array}
\]
\caption{Syntax of \FCdot. In $\ev{pi}\,e\,e'$, in $\obj{\tel}$, and in the telescopes carried by evidence and atoms, the second component lives under one more binder.}
\label{fig:target-syntax}
\end{figure*}

Figure~\ref{fig:target-syntax} gives the syntax. As in the source, everything is intrinsically scoped, and the two calculi share signatures, variables, renamings, and the label type.

\subsection{Types}

A type is $\bot$, a block name $\sel{x}{\ell}$, a dependent function type $\pit{S}{T}$ whose codomain may mention the parameter's block, or an object type $\obj{\tel}$. A telescope is a list of propositions over a self binder, oldest first. Propositions do not bind, so the telescope is a list rather than a dependent context, and $\tat{\tel}{i}{P}$ says that its $i$-th proposition is $P$. The empty object type asserts nothing about its inhabitants and every type is included in it, so $\top$ is defined as $\obj{\tnil}$. Types never contain terms: the only way a type refers to a binder is through the names of its block. Every label has a name in every block, and $\sel{x}{a}$ for a term label $a$ is the type of the field $a$ of $x$.

\subsection{Evidence}

\begin{figure*}
\small
\begin{mathpar}
\inferrule{ }{\lej{\Gamma}{\ev{refl}\,T}{T}{T}}
\and
\inferrule{\lej{\Gamma}{e}{S}{M} \\ \lej{\Gamma}{f}{M}{T}}{\lej{\Gamma}{\ev{trans}\,e\,f}{S}{T}}
\and
\inferrule{ }{\lej{\Gamma}{\ev{top}\,T}{T}{\top}}
\and
\inferrule{ }{\lej{\Gamma}{\ev{bot}\,T}{\bot}{T}}
\and
\inferrule{\eqj{\Gamma}{\varphi}{S}{T}}{\lej{\Gamma}{\ev{eqToLe}\,\varphi}{S}{T}}
\and
\inferrule{\lej{\Gamma}{e}{S_2}{S_1} \\ \lej{\Gamma, x : S_2}{f}{T_1}{T_2}}{\lej{\Gamma}{\ev{pi}\,e\,f}{\pit{S_1}{T_1}}{\pit{S_2}{T_2}}}
\and
\inferrule{\morj{\Gamma}{m}{\tel}{\tel'}}{\lej{\Gamma}{\ev{obj}\,\tel\,m}{\obj{\tel}}{\obj{\tel'}}}
\and
\inferrule{\lej{\Gamma}{e}{S}{\obj{\tel_1}} \\ \lej{\Gamma}{f}{S}{\obj{\tel_2}}}{\lej{\Gamma}{\ev{pair}\,\tel_1\,\tel_2\,e\,f}{S}{\obj{(\tel_1 \mathbin{+\!\!+} \tel_2)}}}
\and
\inferrule{\atomj{\Gamma}{a}{S} \\ \lej{\Gamma}{e}{S}{\obj{\tel}} \\ \tat{\tel}{i}{\ple{S'}{T'}}}
  {\lej{\Gamma}{\ev{member}\,a\,e\,i}{\inst{S'}{\rt{a}}}{\inst{T'}{\rt{a}}}}
\\
\inferrule{\Gamma(x) \ni \ell \mathrel{:=} W}{\eqj{\Gamma}{\ev{def}\,x\,\ell}{\sel{x}{\ell}}{W}}
\and
\inferrule{\atomj{\Gamma}{a}{S} \\ \lej{\Gamma}{e}{S}{\obj{\tel}} \\ \tat{\tel}{i}{\peq{S'}{T'}}}
  {\eqj{\Gamma}{\ev{member}\,a\,e\,i}{\inst{S'}{\rt{a}}}{\inst{T'}{\rt{a}}}}
\and
\inferrule{\atomj{\Gamma}{a}{S} \\ \lej{\Gamma}{e}{S}{\obj{\tel}} \\ \tat{\tel}{i}{\phas{\ell}}}
  {\hasj{\Gamma}{\ev{member}\,a\,e\,i}{\rt{a}}{\ell}}
\and
\inferrule{\Gamma(x) \text{ has fields } \vec{\ell} \\ \ell \in \vec{\ell}}{\hasj{\Gamma}{\ev{field}\,\ell}{x}{\ell}}
\\
\inferrule{ }{\Gamma \vdash \ev{none} : X \le X}
\and
\inferrule{\lej{\Gamma}{e}{A}{B}}{\Gamma \vdash \ev{some}\,e : \weak{A} \le \weak{B}}
\and
\inferrule{ }{\morj{\Gamma}{\ev{nil}}{\tel}{\tnil}}
\and
\inferrule{\morj{\Gamma}{m}{\tel}{\tel'} \\ \holej{\tel}{H}{X}{Y} \\ \Gamma \vdash q : S \le X \\ \Gamma \vdash q' : Y \le T}
  {\morj{\Gamma}{\ev{le}\,m\,q\,H\,q'}{\tel}{\tcons{\tel'}{\ple{S}{T}}}}
\and
\inferrule{\morj{\Gamma}{m}{\tel}{\tel'} \\ \tat{\tel}{j}{\peq{X}{Y}}}
  {\morj{\Gamma}{\ev{eq}\,m\,j\,\mathit{false}}{\tel}{\tcons{\tel'}{\peq{X}{Y}}}}
\and
\inferrule{\morj{\Gamma}{m}{\tel}{\tel'} \\ \tat{\tel}{j}{\peq{X}{Y}}}
  {\morj{\Gamma}{\ev{eq}\,m\,j\,\mathit{true}}{\tel}{\tcons{\tel'}{\peq{Y}{X}}}}
\and
\inferrule{\morj{\Gamma}{m}{\tel}{\tel'} \\ \tat{\tel}{j}{\phas{\ell}}}
  {\morj{\Gamma}{\ev{has}\,m\,j}{\tel}{\tcons{\tel'}{\phas{\ell}}}}
\end{mathpar}
\caption{Evidence typing. The rules for $\ev{refl}$, $\ev{symm}$, and $\ev{trans}$ on equalities are omitted. A hole reads a source proposition as an inclusion: $\holej{\tel}{\ev{le}\,j}{X}{Y}$ when $\tat{\tel}{j}{\ple{X}{Y}}$, $\holej{\tel}{\ev{eq}\,j}{X}{Y}$ when $\tat{\tel}{j}{\peq{X}{Y}}$, and $\holej{\tel}{\ev{eqSym}\,j}{X}{Y}$ when $\tat{\tel}{j}{\peq{Y}{X}}$.}
\label{fig:evidence}
\end{figure*}

There are three families of evidence: inclusions $\lej{\Gamma}{e}{S}{T}$, equalities $\eqj{\Gamma}{\varphi}{S}{T}$, and presence $\hasj{\Gamma}{h}{x}{\ell}$, which says that the object bound at $x$ has a field $\ell$. Figure~\ref{fig:evidence} gives their typing. Endpoints are assigned by typing, not carried by the syntax, and there is no symmetry for inclusions.

The elimination rule $\ev{member}$ is the one rule that connects the type of an atom to the block of its root. From $a : S$ and evidence $e$ for $S \le \obj{\tel}$, the $i$-th proposition of $\tel$ holds at $\rt{a}$, as an inclusion, an equality, or a presence, according to its kind. The rule $\ev{def}$ reads the definition of a block name off a transparent binder. Contexts have two kinds of binder: an opaque one, $x : T$, whose block is abstract, and a transparent one, $x : T \mathrel{:=} W ; \vec\ell$, whose block is defined by the witnesses $W$ and which is known to have the fields $\vec\ell$. A witness list assigns $\top$ to every label it does not mention. Transparent binders arise in two places: inside an object literal, where the fields are typed under the literal's own self binder, and in store typing, where every allocated binder is transparent.

An object coercion $\ev{obj}\,\tel\,m$ is annotated with its source telescope and carries a morphism $m$ that proves each proposition of the target telescope, oldest first. A target inclusion is proven by a template $\ev{le}\,m\,q\,H\,q'$: a hole $H$ names a proposition of the source and reads it as an inclusion $X \le Y$, and the two sides $q$ and $q'$ are closed coercions that adjust the endpoints. A side is either absent or a coercion $A \le B$ between closed types weakened under the self binder, so it cannot mention the self block. A target equality is a source equality, possibly flipped. A target presence is inherited from a source presence by index. A morphism therefore never eliminates through the self block, which is the property Section~\ref{sec:why-structural} relies on. Pairing, $\ev{pair}$, combines two coercions into object types into one coercion into the concatenation, and it is how the source's \rn{And} rule is translated.

\subsection{Atoms, terms, and values}

\begin{figure*}
\small
\begin{mathpar}
\inferrule{ }{\atomj{\Gamma}{x}{\Gamma(x)}}
\and
\inferrule{\atomj{\Gamma}{a}{S} \\ \lej{\Gamma}{e}{S}{T}}{\atomj{\Gamma}{\ev{cast}\,a\,e}{T}}
\and
\inferrule{\atomj{\Gamma}{a}{\obj{\tel}}}{\atomj{\Gamma}{\ev{unfold}\,a}{\obj{\weak{(\inst{\tel}{\rt{a}})}}}}
\and
\inferrule{\atomj{\Gamma}{a}{\obj{\weak{(\inst{\tel}{\rt{a}})}}}}{\atomj{\Gamma}{\ev{fold}\,\tel\,a}{\obj{\tel}}}
\and
\inferrule{\atomj{\Gamma}{a}{\obj{\tel_1}} \\ \atomj{\Gamma}{b}{\obj{\tel_2}} \\ \rt{b} = \rt{a}}{\atomj{\Gamma}{\ev{both}\,\tel_1\,\tel_2\,a\,b}{\obj{(\tel_1 \mathbin{+\!\!+} \tel_2)}}}
\\
\inferrule{\atomj{\Gamma}{a}{T}}{\tmj{\Gamma}{a}{T}}
\and
\inferrule{\valj{\Gamma}{v}{T}}{\tmj{\Gamma}{v}{T}}
\and
\inferrule{\atomj{\Gamma}{a}{\pit{S}{T}} \\ \atomj{\Gamma}{b}{S}}{\tmj{\Gamma}{a\,b}{\inst{T}{\rt{b}}}}
\and
\inferrule{\atomj{\Gamma}{a}{S} \\ \hasj{\Gamma}{h}{\rt{a}}{\ell}}{\tmj{\Gamma}{\ev{proj}\,a\,\ell\,h}{\sel{\rt{a}}{\ell}}}
\and
\inferrule{\tmj{\Gamma}{t}{T} \\ \tmj{\Gamma, x : T}{u}{\weak{U}}}{\tmj{\Gamma}{\klet x = t \kin u}{U}}
\and
\inferrule{\tmj{\Gamma}{t}{S} \\ \lej{\Gamma}{e}{S}{T}}{\tmj{\Gamma}{\ev{cast}\,t\,e}{T}}
\\
\inferrule{\tmj{\Gamma, x : S}{t}{T}}{\valj{\Gamma}{\lambda(x : S).\,t}{\pit{S}{T}}}
\and
\inferrule{\fldj{\Gamma, x : \obj{\ofLit{W}{F}} \mathrel{:=} W ; \dom F}{F}}{\valj{\Gamma}{\lit{W}{F}}{\obj{\ofLit{W}{F}}}}
\and
\inferrule{\valj{\Gamma}{v}{S} \\ \lej{\Gamma}{e}{S}{T}}{\valj{\Gamma}{\ev{cast}\,v\,e}{T}}
\and
\inferrule{ }{\fldj{\Gamma}{\tnil}}
\and
\inferrule{\fldj{\Gamma}{F} \\ \tmj{\Gamma}{t}{\sel{x}{\ell}}}{\fldj{\Gamma}{F, \ell = t}}
\end{mathpar}
\caption{Typing of atoms, terms, values, and fields. In the field rules $x$ is the self binder of the literal.}
\label{fig:terms}
\end{figure*}

An atom is a variable under wrappers that erase to nothing: a cast, the unfolding of a self block at the root, the folding of the root back into a self block, and the pairing of two typings of one variable. Figure~\ref{fig:terms} gives the typing. Unfolding instantiates the self binder of the telescope by the atom's root and weakens the result under a fresh, unused self binder, so that $\obj{\weak{(\inst{\tel}{\rt{a}})}}$ is an object type whose propositions are about $\rt{a}$ rather than about the self. Folding is its inverse. These are the atom forms of the source rules \rn{Rec-E} and \rn{Rec-I}, and $\ev{both}$ is \rn{And-I}.

Terms are in monadic normal form over atoms: application and projection take atoms, a let sequences computations, and a cast on a term becomes a cast frame when it runs. There is no subsumption. Application instantiates the codomain at the root of the argument, as \rn{All-E} does in the source, and the result of a let may not mention the block of the let-bound variable. A projection $\ev{proj}\,a\,\ell\,h$ carries the presence evidence $h$ and has the block name $\sel{\rt{a}}{\ell}$ as its type. Converting that name to a declared field type takes a further cast, which the translation supplies from the declaration.

An object literal $\lit{W}{F}$ has witnesses and fields and nothing else. Its type is its precise type $\obj{\ofLit{W}{F}}$: the telescope with one equality $\peq{\sel{\self}{\ell}}{W(\ell)}$ per witness, followed by one presence $\phas{\ell}$ per field. Fields are typed under the self binder bound transparently at that type, and each field body has the block name $\sel{x}{\ell}$ of its own label as its type. A literal typed at a declaration type is a literal under a cast. This is the inert-store design of Section~\ref{sec:why-structural}, and it is what makes a witness of a literal free to be a name of the same block: a definition $\ell = \sel{x}{\ell'}$ is an alias, and Section~\ref{sec:resolution} explains how aliases, including cyclic ones, are resolved.

\subsection{Stores and the machine}

\begin{figure*}
\small
\[
\begin{array}{@{}lcll@{}}
\mst{\sigma}{K}{\klet x = t \kin u} &\step& \mst{\sigma}{K \triangleright \ev{let}\,u}{t} \\
\mst{\sigma}{K}{\ev{cast}\,t\,e} &\step& \mst{\sigma}{K \triangleright \ev{cast}\,e}{t} \\
\mst{\sigma}{K \triangleright \ev{cast}\,e}{v} &\step& \mst{\sigma}{K}{\ev{cast}\,v\,e} \\
\mst{\sigma}{K \triangleright \ev{cast}\,e}{a} &\step& \mst{\sigma}{K}{\ev{cast}\,a\,e} \\
\mst{\sigma}{K \triangleright \ev{let}\,u}{v} &\step& \mst{\sigma, \mathit{core}(v)}{\weak{K}}{\mathit{adjust}(u, v)} \\
\mst{\sigma}{K \triangleright \ev{let}\,u}{a} &\step& \mst{\sigma}{K}{\inst{u}{a}} \\
\mst{\sigma}{K}{x\,b} &\step& \mst{\sigma}{K}{\inst{t_0}{b}} & \text{if } \sigma(x) = \lambda(z : S_0).\,t_0 \\
\mst{\sigma}{K}{a\,b} &\step& \mst{\sigma}{K}{\inst{t_0}{b}} & \text{if } \sigma(\rt{a}) = \lambda(z : S_0).\,t_0,\ a \neq \rt{a},\ \chn{\sigma}{a}{n}{(a', F)},\ F \in \{\form{id}, \form{eqv}\,\varphi\} \\
\mst{\sigma}{K}{a\,b} &\step& \mst{\sigma}{K}{\ev{cast}\,(\inst{t_0}{\ev{cast}\,b\,d})\,(\inst{c}{b})} & \text{if } \sigma(\rt{a}) = \lambda(z : S_0).\,t_0,\ a \neq \rt{a},\ \chn{\sigma}{a}{n}{(a', \form{pi}\,d\,c)} \\
\mst{\sigma}{K}{\ev{proj}\,a\,\ell\,h} &\step& \mst{\sigma}{K}{\inst{t}{\rt{a}}} & \text{if } \sigma(\rt{a}) = \lit{W}{F} \text{ and } F(\ell) = t
\end{array}
\]
\caption{The store machine of \FCdot. $\mathit{core}(v)$ is the literal under the cast wrappers of $v$, and $\mathit{adjust}(u, v)$ substitutes $\ev{cast}\,x\,E$ for $x$ in $u$ when $E$ is the composite of those wrappers, and is $u$ when there are none. The judgment $\chn{\sigma}{a}{n}{(a', F)}$ is the normalization of the chain of casts of $a$ with fuel $n$, defined in Section~\ref{sec:normalizer}.}
\label{fig:machine}
\end{figure*}

A store holds one value per binder of its signature, and store typing $\storej{\sigma}{\Gamma}$ requires every entry to be a literal, typed in the transparent context of the entries before it, with the context recording the literal's witnesses and field labels:
\begin{mathpar}
\inferrule{ }{\storej{\cdot}{\cdot}}
\and
\inferrule{\storej{\sigma}{\Gamma} \\ v \text{ is a literal} \\ \valj{\Gamma}{v}{T}}{\storej{\sigma, v}{\Gamma, x : T \mathrel{:=} \mathit{witnesses}(v) ; \mathit{fields}(v)}}
\end{mathpar}
A lambda has no witnesses and no fields, so its block names are all defined as $\top$. A continuation is a list of frames, a let frame $\ev{let}\,u$ or a cast frame $\ev{cast}\,e$, and $\contj{\Gamma}{K}{T}{U}$ says that $K$ accepts a value of type $T$ and produces one of type $U$. A state is typed at $U$ when its store is typed in a context in which the term has some type $T$ and the continuation takes $T$ to $U$.

Figure~\ref{fig:machine} gives the steps. Allocation is the rule that keeps stores inert: the value that reaches a let frame may carry casts, and the machine stores the literal underneath at its own type and rewrites the continuation so that the new variable is used under the composite of the stripped casts. Application through a bare variable is ordinary $\beta$. Application through a wrapped atom reads the domain and codomain evidence off the normal form of the atom's chain of casts: when that form is the identity or a conversion, the closure is applied directly, and when it is a function coercion $\form{pi}\,d\,c$, the argument is cast by $d$ and the result by $c$ instantiated at the argument. The machine thus runs the normalizer of Section~\ref{sec:canonical}, and progress has to show that the normalizer succeeds on a closed atom of function type. Projection reads the field off the literal at the root and instantiates the field's self binder by the root.

Atom substitution $\inst{u}{a}$ replaces the variable by the atom in terms, and renames the variable to $\rt{a}$ in types and evidence, which see nothing of an atom but its root.

\subsection{Erasure}

Erasure into the runtime maps an atom to its root, drops every cast and every piece of evidence, keeps the fields of a literal, and drops cast frames. A step of the machine either moves a cast frame, in which case the erased state does not change, or erases to exactly one runtime step. Cast-frame steps are bounded by the nesting of casts in the state, so they cannot be taken forever.

\subsection{The checker}

Every constructor of evidence and of atoms determines its endpoints from its arguments and the context: object coercions carry their source telescope, pairing and folding carry their target telescopes, and a template is checked from its hole outwards. The checker is therefore a family of synthesizing kernels, one per judgment, each returning the typing derivation it found. Soundness holds by construction, and completeness, that every derivation is found, is a theorem proven by induction on derivations. Its top-level form is $\lean{checkTm}\ \Gamma\ t\ T = \mathit{true} \iff \tmj{\Gamma}{t}{T}$, and all examples of Section~\ref{sec:mechanization} are checked by running it inside the kernel of the proof assistant.
